\documentclass[12pt]{article}
\usepackage[T1]{fontenc}
\usepackage[utf8]{inputenc}
\usepackage{lmodern}
\usepackage{textcomp}
\usepackage{enumitem}
\usepackage{geometry}
\geometry{margin=1in}
\begin{document}
\begin{center}
Answer Key - Computer Science - Graduate Level Assessment 3 \
\small (Answers are ordered exactly as in the question paper.)
\end{center}

\section*{Multiple Choice}
\begin{enumerate}[label=\arabic*.]
\item \textbf{Answer:} B
\\ \textit{Options:} A) Integrity; B) Condentiality; C) Authentication; D) Nonrepudiation
\\ \textit{Note:} The correct answer is "Confidentiality" because message confidentiality ensures that the content of a communication remains private and is accessible only to the intended sender and receiver, thereby protecting sensitive information from unauthorized access.
\item \textbf{Answer:} A
\\ \textit{Options:} A) WEP; B) WPA; C) WPA 2; D) WPA 3
\\ \textit{Note:} WEP (Wired Equivalent Privacy) is considered the least strong security encryption standard because it uses outdated and weak encryption algorithms, making it vulnerable to various attacks and easily compromised compared to more modern standards like WPA and WPA 2.
\item \textbf{Answer:} D
\\ \textit{Options:} A) Haunted web; B) World Wide Web; C) Surface web; D) Deep Web
\\ \textit{Note:} The Deep Web refers to parts of the internet that are not indexed by standard search engines, making them inaccessible through typical search queries.
\item \textbf{Answer:} C
\\ \textit{Options:} A) Insertion sort; B) Quicksort; C) Merge sort; D) Selection sort
\\ \textit{Note:} Merge sort consistently divides the input array into smaller subarrays and merges them in a sorted manner, resulting in a time complexity of O(n log n) regardless of the initial order of the elements, making its performance least affected by input arrangement compared to other algorithms.
\item \textbf{Answer:} A
\\ \textit{Options:} A) IM - Trojans; B) Backdoor Trojans; C) Trojan-Downloader; D) Ransom Trojan
\\ \textit{Note:} IM Trojans are specifically designed malware that infiltrate instant messaging applications to capture and transmit users' login credentials and passwords.
\end{enumerate}

\section*{True / False}
\begin{enumerate}[label=\arabic*.]
\setcounter{enumi}{5}
\item \textbf{Answer:} True
\\ \textit{Options:} A) True; B) False
\\ \textit{Note:} Escaping queries is true because it sanitizes user input by encoding potentially harmful characters, thereby preventing them from being executed as part of SQL commands, which effectively mitigates the risk of SQL injection attacks.
\item \textbf{Answer:} False
\\ \textit{Options:} A) True; B) False
\\ \textit{Note:} The statement is false because a stack is not a sequential segment of memory; rather, it is a data structure that follows the Last In, First Out (LIFO) principle for managing function calls and local variables, distinct from arrays or character strings.
\item \textbf{Answer:} True
\\ \textit{Options:} A) True; B) False
\\ \textit{Note:} Finding a shortest cycle in an undirected graph is solvable in polynomial time because it can be efficiently determined using algorithms such as breadth-first search (BFS) which explores all possible paths and identifies the shortest cycle.
\item \textbf{Answer:} False
\\ \textit{Options:} A) True; B) False
\\ \textit{Note:} A packet filter firewall operates at the network and transport layers, analyzing header information rather than inspecting application data for security threats.
\item \textbf{Answer:} True
\\ \textit{Options:} A) True; B) False
\\ \textit{Note:} The statement is true because Confidentiality, Integrity, Non-repudiation, and Authentication are fundamental principles that ensure secure communication and data protection in information security.
\end{enumerate}

\section*{Long Form Response}
\begin{enumerate}[label=\arabic*.]
\setcounter{enumi}{10}
\item \textbf{Answer:} def intersection\_nested\_lists(l 1, l 2): | return result
\\ \textit{Note:} correct because it defines a function intended to compute the intersection of elements from a nested list (l 1) with another list (l 2), signifying the need for an algorithm that identifies common elements between these two data structures.
\item \textbf{Answer:} def max\_Abs\_Diff(arr,n): | return (maxEle - minEle)
\\ \textit{Note:} The answer correctly identifies that the maximum difference between any two elements in an array can be determined by subtracting the minimum element from the maximum element, effectively capturing the largest possible gap between values in the dataset.
\end{enumerate}

\vfill
\noindent\textit{End of Answer Key}
\end{document}